% Generated draft: requires mathematical and source review.
\documentclass[11pt]{article}
\usepackage[margin=25mm]{geometry}
\usepackage{fontspec}
\setmainfont{Latin Modern Roman}
\usepackage{amsmath,amssymb,mathtools,unicode-math}
\setmathfont{Latin Modern Math}
\usepackage{enumitem,needspace,xurl,hyperref}
\setlength{\parindent}{0pt}
\setlength{\parskip}{5pt}
\setlength{\emergencystretch}{3em}
\newcommand{\N}{\mathbb{N}}
\newcommand{\Z}{\mathbb{Z}}
\newcommand{\Q}{\mathbb{Q}}
\newcommand{\R}{\mathbb{R}}
\newcommand{\C}{\mathbb{C}}
\newcommand{\F}{\mathbb{F}}
\newcommand{\A}{\mathbb{A}}
\newcommand{\PP}{\mathbb{P}}
\newcommand{\E}{\mathbb{E}}
\newcommand{\eps}{\varepsilon}
\newcommand{\dd}{\,\mathrm{d}}
\begin{document}
\section*{UnsolvedMath ID: 20000805}
\textbf{Source problem code:} catalogue entry
\subsection{Definitions and mathematical statement}
The source asks two existence questions: whether a component of the Hilbert scheme of $d$ points in affine $n$-space can exist only in characteristic $p$, for some prime $p$, and the analogous question for Hilbert schemes of curves in $\PP^3$. To make the aid's strongest interpretation precise, use the relative schemes over $\Z$. A $p$-vertical component is an irreducible component of the total relative Hilbert scheme supported entirely over the closed point $(p)$; a horizontal component dominates $\operatorname{Spec}\Z$. A component merely appearing in a special fiber need not be vertical in this sense.

The supplied aid attributes to Jelisiejew, Corollary 5.2, the affirmative affine statement
\[
 \text{for every prime }p\text{ there exists }d\text{ such that }
 \operatorname{Hilb}^{d}(\A^{16}_{\Z})
 \text{ has a }p\text{-vertical irreducible component}.
\]
The length $d$ may depend on $p$. The record gives no construction or full bibliographic identification of that corollary.

For the projective question, fix a Hilbert polynomial $P$ and write $\mathcal H=\operatorname{Hilb}^{P}(\PP^n_{\Z})$. The aid reports a necessary boundary-incidence condition: every hypothetical $p$-vertical irreducible component, after passage to the appropriate geometric special fiber, contains a Borel-fixed monomial point $[I]$. The monomial ideal lifts by the same monomials over $\Z$, giving a flat integral family with that Hilbert polynomial. Consequently $[I]$ also lies on a horizontal component. In characteristic $p$, Borel-fixed must not be silently replaced by characteristic-zero strong stability.

Equivalently, at such a boundary point the relevant local relative Hilbert ring has distinct minimal primes, one defining a vertical branch over $p$ and one defining a horizontal branch. Thus the proposed search for vertical components of $\operatorname{Hilb}^{P}(\PP^3_{\Z})$ can examine mixed vertical/horizontal minimal primes at integrally liftable Borel-fixed monomial ideals. The existence of the integral monomial lift does not remove the possible vertical branch, and the reduction does not itself exhibit one.

The aid says it found no answer for curves in fixed $\PP^3$; this is its search assessment, not an independently verified current literature claim. The polynomial $P$, degree/genus and meaning of ``curves'' remain unspecified. The proper projective Hilbert scheme includes schemes with embedded points; if one instead imposes purity, smoothness or local Cohen--Macaulay conditions, the Borel-fixed degeneration need not remain inside that restricted open locus. These additional restrictions are not supplied by the source.
\subsection{Short English statement}
The aid reports a positive answer for affine points: every prime supports a vertical component for some length in affine sixteen-space. The fixed-$\PP^3$ curve question is left unresolved. Its boundary reduction says a hypothetical vertical component must meet a horizontal one at a liftable monomial point, so liftability of that point alone is no contradiction.
\subsection{Sources}
\begin{itemize}
\item Local UnsolvedMath registry, permanent ID $20000805$, source code \texttt{catalogue entry}, supplied complete source record and research summary.
\item American Institute of Mathematics, \emph{Components of Hilbert Schemes}, Problem 4, supplied URL \url{https://aimath.org/WWN/hilbertschemes/hilbertschemes.pdf}; canonical local record \texttt{aim-arithmetic-geometry-notes.json}, \texttt{notes[50]}. No PDF was retrieved.
\item Jelisiejew, Corollary 5.2, as cited in the supplied summary for the affine-points assertion. No title, year, identifier or full bibliographic reference is supplied. The original literature field is empty.
\item Research attempt 1: partial result; problem partially solved at the recorded run; supplied rigor label ``proved'', novelty confidence ``low''. The machine-generated aid is not peer reviewed and requires independent verification before citation; its remaining-open assessment was not independently updated.
\end{itemize}
\end{document}
